\section{Further research questions}
\label{cont:sec:questions}

The following programme prioritizes questions for which the present
results supply a concrete starting point. Some are new refinements
of the baseline questions; others remain difficult after the
progress made here. The statements in this section are proposals,
not additional theorems.

\subsection{Harmonic sums and uniform asymptotics}

\begin{enumerate}[leftmargin=*,label=H\arabic*.]
\item \textbf{Certified relative remainders on the entire saddle
range.} The proof of Theorem~\ref{us:thm:uniform} supplies uniform
derivative and curvature control, including the explicit bound
\eqref{us:eq:curvature-bound}. Derive useful computable constants
$C_J,p_J$ such that
\[
 \left|\frac{R_{p,h}}{\mathcal S_{p,h}}
      -\sum_{j=0}^JE_jp^{-j}\right|
 \le C_Jp^{-J-1}\quad(p\ge p_J)
\]
holds for every $h\ge1,a>0$. The main task is to retain practical
constants in the Gaussian Taylor and endpoint bounds.

\item \textbf{A bridge to bounded $p$.} The new theorem is uniform
in $p/h$, but requires $p\to\infty$. Find a single approximation
which joins it to the fixed-$p$, growing-$h$ endpoint expansion.
An incomplete-gamma or gamma-density formulation may be more
natural than forcing a Gaussian approximation when its shape
parameter stays bounded.

\item \textbf{Inverse fractions approaching zero or one.}
The inverse theorem keeps $\theta$ away from both endpoints.
Determine the largest moving ranges of
$\lambda=-\log\theta$ for which the polynomial inverse remains
uniform. For example, how do the scales change if
$\lambda\to0$ or $\lambda\to\infty$ with $h$? Near the lower
endpoint $(\log2)^h/2$, the transition underlying the inverse
itself changes.

\item \textbf{Late saddle coefficients.} The all-orders expansion
is effective at every fixed order, but does not describe the growth
of $E_j$ as $j\to\infty$. Identify the complex saddles and
singularities controlling that growth, and determine whether an
optimal-truncation theorem can be uniform over the moving real
saddle. This would connect the harmonic analysis with the
manuscript's sharp reflected-moment work.
\end{enumerate}

\subsection{Two growing log-gamma exponents}

\begin{enumerate}[leftmargin=*,label=M\arabic*.]
\item \textbf{An analytic slope-ratio inequality and its sharp
constant.} Theorem~\ref{gs:thm:monotonicity} proves global
monotonicity using an exact interval certificate in the middle.
Find a proof based entirely on analytic inequalities, and determine
the optimal lower bound for $J'/J$. The symmetry suggests that
the midpoint deserves special attention, but the certificate given
here does not assert that it is the minimizer.

\item \textbf{Moving transition parameter.} Extend
Theorem~\ref{joint:thm:transition} to transition parameters which
tend to zero or infinity. Its correction polynomials exhibit
how powers of $\lambda$ compete with powers of $m^{-1}$, but
that formal observation is not a uniform bound. Determine the
first range in which the cubic and higher local perturbations
change the leading transition function.

\item \textbf{Uniform bridges between moment regimes.}
Construct an approximation valid while the dominant saddle
travels from a fixed interior point to $x\asymp m^{-1}$ and
beyond. The transition limit proved here supplies a necessary
boundary condition. A successful formula must also control the
opposite endpoint, rather than using a local series alone.

\item \textbf{A general analytic endpoint-pair theorem.}
For $f(x)=-\log x+g(x)$ and
$B(x)=b_1x+b_2x^2+\cdots$, the local argument suggests the
leading factor
\[
 \exp\bigl((g_1+b_2/b_1)\lambda\bigr).
\]
State and prove explicit global conditions under which this
generalization is valid. The exact coefficient recipe indicates
which parts are formal; the endpoint localization identifies
where additional hypotheses are needed.

\item \textbf{Optimal truncation with a growing reflected
exponent.} Determine the maximal growth of $m$ for which the
fixed-$m$ half-first-omitted-term law remains uniform. The factor
$B(-\tau)^m$ in the inverse-gamma amplitude changes the late-term
saddle. A theorem must track this effect together with the
transition polynomials, not simply substitute a variable $m$
into a fixed-amplitude estimate.
\end{enumerate}

\subsection{Integral distribution lattices}

\begin{enumerate}[leftmargin=*,label=D\arabic*.]
\item \textbf{Three or more prime divisors.} Determine
$\operatorname{coker}P_q$ integrally for squarefree $q$ with at
least three prime divisors. The two-prime core suggests an
incidence presentation involving intersections of divisor layers.
Which torsion classes require more than pairwise cycle information?

\item \textbf{Prime powers in the primitive cokernel.} Extend
the explicit classification from $p\ell$ to $p^a\ell^b$.
The full quotient already has a basis and finite resolution, so
the difficulty is the primitive inclusion. Separate the effect
of higher primary heights from multiplication-cycle torsion.

\item \textbf{An exact denominator formula for every level.}
For $q=p\ell$ the answer is a least common multiple of three
explicit factors. For general $q$, seek a formula for the exponent
of the integral primitive cokernel, rather than merely its
determinant. A product bound can exceed the needed denominator
substantially, as the level-fifteen example demonstrates.

\item \textbf{Integral reflection and jet modules.} Study
reflection quotients in characteristic two and spectral jets over
integer or finite-characteristic coefficient rings. The full
distribution module remains free, but the primitive image may
simultaneously have spectral defects and modular torsion. A
presentation that keeps both effects should be useful for
special-value computations without semisimple character projectors.
\end{enumerate}

\subsection{Gaussian polylogarithms and exact discovery}

\begin{enumerate}[leftmargin=*,label=G\arabic*.]
\item \textbf{Prove or refute the common $S_6$ conjecture.}
Use the reduced coordinate form
\eqref{s6audit:eq:reduced-candidate}. The weight-five success
indicates that transformations of iterated-integral paths can
supply information outside a restricted depth-two product module.
A successful proof must display the extra analytic relation and
an exact finite combination, not only a numerical null vector.

\item \textbf{Control relation-space growth before elimination.}
The weight-seven experiment shows how sparse rows can become
dense during elimination. Work first in shuffle indecomposables,
exploit color and reversal symmetries, and isolate the relevant
quotient before adjoining transformed paths. A modular calculation
which suggests membership should be followed by rational
reconstruction and exact verification of every relation.

\item \textbf{Obstructions with an explicit relation vocabulary.}
Extend the existing uniform depth-two obstructions to carefully
specified larger presentations. A nonzero functional on a
restricted module is useful because it identifies missing
information; it is not a theorem about all functional equations
of the evaluated periods. State the alphabet, regularization
rules, allowed depths, and parameter transformations in each
obstruction theorem.

\item \textbf{A reproducible conjecture registry.} Store each
candidate with an immutable coordinate order, a primitive integer
relation, discovery precision, frozen-coefficient reevaluation,
and its exact relations to other candidates. The two $S_6$
presentations show why coordinate equivalence should be checked
before counting independent discoveries. Keep numerical evidence,
formal proofs, and interval certificates as distinct fields.
\end{enumerate}

These questions preserve the central aim of the programme: explain
why an identity holds, identify the structure that makes it
computable, and state precisely where an argument ceases to be
uniform or a formal relation ceases to determine an arithmetic one.
